\documentclass[10pt,a4paper]{article}
\usepackage[hmargin=1.5cm,top=1.1cm,bottom=1.1cm]{geometry}
\usepackage[T1]{fontenc}
\usepackage[utf8]{inputenc}
\usepackage{lmodern}
\usepackage{microtype}
\usepackage{booktabs}
\usepackage{graphicx}
\usepackage{xcolor}
\usepackage{tcolorbox}
\usepackage{enumitem}
\usepackage[hidelinks]{hyperref}
\usepackage{soul}
\sethlcolor{yellow!55}
\soulregister\texttt7
\soulregister\emph7
\soulregister\textbf7

\setlist{nosep,leftmargin=*}
\setlength{\parskip}{2.5pt}
\setlength{\parindent}{0pt}
\pagestyle{empty}
\newcommand{\figdir}{../output/figures}
\newcommand{\q}[1]{\textbf{\color{blue!50!black}#1}}

\begin{document}
\small

{\large\textbf{Predicting employee attrition with decision trees}}\\
Ignacio Queipo de Llano \textperiodcentered{} Predictive Machine Learning in Finance, IESE MiF \textperiodcentered{} 12 November 2026\\[-6pt]
\rule{\linewidth}{0.4pt}

\begin{tcolorbox}[colback=white,colframe=black,boxrule=0.4pt,sharp corners,left=5pt,right=5pt,top=3pt,bottom=3pt]
\textbf{Bottom line.} Only 16.1\% of the 1,470 employees left, so a model that always says ``Stay'' is already 84\%
accurate. The best tree (depth 5) reaches 84.7\% on the test set and 84.1\% in cross-validation: \emph{in accuracy it
barely beats doing nothing}. Its value is elsewhere: it finds a third of the leavers with few false alarms and turns
them into rules a manager can read. Logistic regression is better on every metric (87.8\% accuracy, 40\% of leavers
caught, AUC 0.83 vs 0.67). Recommendation: score with the logit, lower the alarm to 25\% (catches 68\% of leavers),
and use a depth-3 tree to explain \emph{who} is at risk.
\end{tcolorbox}

\q{Q1. Attrition rate:} \hl{16.1\% (237 of 1,470).} \q{Categorical vs numeric:} \hl{7 text columns are categorical}
(\texttt{Department}, \texttt{JobRole}, \texttt{MaritalStatus}, \texttt{BusinessTravel}, \texttt{EducationField},
\texttt{Gender}, \texttt{OverTime}); the 26 integer columns look numeric. \q{Strongest visual relationship} (attrition
rate by group; leaver vs stayer means): \hl{\texttt{JobRole} (Sales Reps 40\% vs Research Directors 2.5\%),
\texttt{OverTime} (31\% vs 10\%), \texttt{JobLevel} 1 (26\%), \texttt{Single} (26\%), low \texttt{JobInvolvement}
(34\%); leavers are younger, earlier in their career and paid less} (about 0.45 s.d. below stayers).

\q{Q2. Types:} \hl{nominal text, binary text, ordinal integer codes (satisfaction 1--4, \texttt{Education},
\texttt{JobLevel}, \texttt{StockOptionLevel}), true numbers (14), and 4 useless columns} (3 constant, 1 ID: drop).
\q{Hidden?} \hl{Yes: the ordinal codes are categories with an order, not quantities.} \q{Treatment:} \hl{true numbers and
ordinal codes as they are; binary as 0/1; nominal one-hot.} \q{Does a tree care?} \hl{Barely: it only asks whether
$x \le$ a threshold, so only the order matters} (a logit would treat 4 as twice 2). Features: 20 on pay, overtime,
career stage, satisfaction, family and department (24 columns after encoding). \hl{Stratified 80/20 split: 1,176 / 294,
16\% leavers in both.}

\vspace{2pt}
\begin{minipage}[t]{0.40\linewidth}\vspace{0pt}
\includegraphics[width=\linewidth]{\figdir/q3_accuracy_vs_depth.png}
\end{minipage}\hfill
\begin{minipage}[t]{0.57\linewidth}\vspace{0pt}
\q{Q3. Depth.} Train accuracy rises to 100\% with no limit (151 leaves: it memorises). Test accuracy peaks at
\hl{depth 5 (84.7\%) and falls from there: overfitting starts after 5.} But the curve is flat between 2 and 5,
and redoing the split with 50 seeds, depth 2 or 3 wins 34 times and depth 5 only 8: ``5'' is partly luck of this split.

\q{Q4. Cross-validation (depth 5).} Folds 82.3--85.7\%, \hl{mean 84.1\%, s.d. 1.3 pp.} Slightly below
the single test score (84.7\%), which was the lucky draw we picked the depth on. \hl{The s.d. is small, so the model is
stable,} but 1.3 pp is six times its edge over ``always Stay'' (83.9\%, a 0.2 pp gap).

\q{Q6. Importances.} \hl{\texttt{TotalWorkingYears} 0.27, \texttt{OverTime} 0.13, \texttt{MonthlyIncome} 0.12,
\texttt{Age} 0.09, \texttt{Single} 0.06. Matches intuition (early career, overworked, underpaid).} Satisfaction scores
matter little in the tree but a lot in the logit: the tree uses a few strong splits, the logit adds up many small effects.
\end{minipage}

\vspace{3pt}
\q{Q5 and Q8. Confusion matrix and model comparison} (test set, 294 employees, 47 leavers):

\vspace{2pt}
\centerline{\begin{tabular}{lrrrrrrrr}
\toprule
 & TP & FN & FP & TN & Accuracy & TP rate (recall) & FP rate & AUC\\
\midrule
Always ``Stay'' & 0 & 47 & 0 & 247 & 84.0\% & 0\% & 0\% & 0.50\\
Tree, depth 5 & 16 & 31 & 14 & 233 & 84.7\% & \hl{34.0\%} & \hl{5.7\%} & 0.67\\
Logistic regression & 19 & 28 & 8 & 239 & \textbf{87.8\%} & 40.4\% & \textbf{3.2\%} & \textbf{0.83}\\
Logit, alarm at 25\% & 32 & 15 & 36 & 211 & 82.7\% & \textbf{68.1\%} & 14.6\% & 0.83\\
\bottomrule
\end{tabular}}

\vspace{2pt}
\emph{HR reading.} \hl{A \textbf{false negative} is a leaver we missed: replacing them costs 50--200\% of a salary. A
\textbf{false positive} is a stayer we flagged: the cost is a retention conversation or bonus spent on someone who was not
leaving. Missing a leaver is far more expensive,} so accuracy at the default 50\% cut-off is the wrong target: the tree
misses 2 of every 3 leavers. Lowering the alarm trades cheap false positives for expensive false negatives.

\vspace{2pt}
\centerline{\includegraphics[width=0.86\linewidth,trim=0 15 0 25,clip]{\figdir/q7_tree_depth3.png}}

\q{Q7. Business rules (depth-3 tree).} Each box: share of employees, [stay, leave] proportions.
\begin{enumerate}[label=\arabic*.]
\item \hl{\textbf{If} total working years $\le$ 1 \textbf{and} age $\le$ 33 \textbf{and} overtime: \textbf{90\% leave}} (20 people).
\item \hl{\textbf{If} more than 1 year of experience \textbf{and} overtime \textbf{and} monthly income $\le$ \$2,476: \textbf{60\% leave}} (40 people).
\item \hl{\textbf{If} more than 1 year of experience, no overtime and work-life balance above Bad: \textbf{8\% leave}} --- the safe
63\% of the workforce. (One leaf holds a single employee: that is noise, not a rule.)
\end{enumerate}

\q{Q8. Recommendation.} \hl{Use the \textbf{logistic regression} to score everyone: higher accuracy, more leavers caught,
fewer false alarms, and it ranks risk much better (AUC 0.83).} With the alarm at 25\% HR reaches 68\% of leavers. \hl{Use
the \textbf{depth-3 tree} to explain the drivers in a meeting (overtime, first job, low pay);} a tree deep enough to compete
is 24 leaves and no longer readable.
\end{document}
